\subsection{Семинар <<Работа со списками>> (4 часа)}

В данной работе вам необходимо взять код с семинара 
<<Объектно-ориентированное программирование>> и разделить его на три части:
\begin{itemize}
\item модуль \textbf{core} (наведите курсор мыши на проект, нажмите правую кнопку мыши и создайте 
модуль Java/Kotlin library);
\item модуль \textbf{console} (консольное приложение, которое работает как раньше и использует
модуль core);
\item модуль \textbf{app} (он создается мастером по умолчанию; Android-приложение, которое
использует модуль core).
\end{itemize}

В последних двух модулях подключение core делается посредством строки:

\begin{verbatim} 
implementation(project(":core"))
\end{verbatim}

Реализуйте функциональность консольного приложения в Android-приложении, 
не забыв выделить отдельно View, ViewModel, а также Model, которая преимущественно
реализуется модулем core. Используйте Navigation для переходов между экранами.

Дополните Android-приложение полной CRUD-функциональностью (создание, просмотр, обновление
и удаление сущностей согласно вашему варианту).
